\documentclass[10pt,a4paper]{amsart}
\usepackage{amsmath,amssymb,amsthm,graphicx}
\usepackage{verbatim}
\newenvironment{reference}{\comment}{\endcomment}
%\newenvironment{reference}{}{}
\newenvironment{slogan}{\comment}{\endcomment}
\newenvironment{history}{\comment}{\endcomment}

% For commutative diagrams we use Xy-pic
\usepackage[all]{xy}

% We use 2cell for 2-commutative diagrams.
\xyoption{2cell}
% Theorem environments.
%
\theoremstyle{plain}
\newtheorem{theorem}[subsection]{Theorem}
\newtheorem{proposition}[subsection]{Proposition}
\newtheorem{lemma}[subsection]{Lemma}

\theoremstyle{definition}
\newtheorem{definition}[subsection]{Definition}
\newtheorem{example}[subsection]{Example}
\newtheorem{exercise}[subsection]{Exercise}
\newtheorem{situation}[subsection]{Situation}

\theoremstyle{remark}
\newtheorem{remark}[subsection]{Remark}
\newtheorem{remarks}[subsection]{Remarks}

\numberwithin{equation}{subsection}

% Macros
%
\def\lim{\mathop{\mathrm{lim}}\nolimits}
\def\colim{\mathop{\mathrm{colim}}\nolimits}
\def\Spec{\mathop{\mathrm{Spec}}}
\def\Hom{\mathop{\mathrm{Hom}}\nolimits}
\def\Ext{\mathop{\mathrm{Ext}}\nolimits}
\def\SheafHom{\mathop{\mathcal{H}\!\mathit{om}}\nolimits}
\def\SheafExt{\mathop{\mathcal{E}\!\mathit{xt}}\nolimits}
\def\Sch{\mathit{Sch}}
\def\Mor{\mathop{\mathrm{Mor}}\nolimits}
\def\Ob{\mathop{\mathrm{Ob}}\nolimits}
\def\Sh{\mathop{\mathit{Sh}}\nolimits}
\def\NL{\mathop{N\!L}\nolimits}
\def\CH{\mathop{\mathrm{CH}}\nolimits}
\def\proetale{{pro\text{-}\acute{e}tale}}
\def\etale{{\acute{e}tale}}
\def\QCoh{\mathit{QCoh}}
\def\Ker{\mathop{\mathrm{Ker}}}
\def\Im{\mathop{\mathrm{Im}}}
\def\Coker{\mathop{\mathrm{Coker}}}
\def\Coim{\mathop{\mathrm{Coim}}}

% Boxtimes
%
\DeclareMathSymbol{\boxtimes}{\mathbin}{AMSa}{"02}

%
% Macros for moduli stacks/spaces
%
\def\QCohstack{\mathcal{QC}\!\mathit{oh}}
\def\Cohstack{\mathcal{C}\!\mathit{oh}}
\def\Spacesstack{\mathcal{S}\!\mathit{paces}}
\def\Quotfunctor{\mathrm{Quot}}
\def\Hilbfunctor{\mathrm{Hilb}}
\def\Curvesstack{\mathcal{C}\!\mathit{urves}}
\def\Polarizedstack{\mathcal{P}\!\mathit{olarized}}
\def\Complexesstack{\mathcal{C}\!\mathit{omplexes}}
% \Pic is the operator that assigns to X its picard group, usage \Pic(X)
% \Picardstack_{X/B} denotes the Picard stack of X over B
% \Picardfunctor_{X/B} denotes the Picard functor of X over B
\def\Pic{\mathop{\mathrm{Pic}}\nolimits}
\def\Picardstack{\mathcal{P}\!\mathit{ic}}
\def\Picardfunctor{\mathrm{Pic}}
\def\Deformationcategory{\mathcal{D}\!\mathit{ef}}
\usepackage{cleveref}
\crefname{originallemma}{original lemma}{original lemmas}
\Crefname{originallemma}{Original lemma}{Original lemmas}
\makeatletter
\newlabel{descent-lemma-closed-immersion}{{03I0}{source}{}{sourcepointer-1}{}}
\newlabel{descent-lemma-closed-immersion@cref}{{[originallemma][0][]03I0}{[1][0][]source}{}{}{}}
\makeatother
\crefname{originallemma}{original lemma}{original lemmas}
\Crefname{originallemma}{Original lemma}{Original lemmas}
\makeatletter
\newlabel{descent-lemma-descent-data-sheaves}{{02W5}{source}{}{sourcepointer-2}{}}
\newlabel{descent-lemma-descent-data-sheaves@cref}{{[originallemma][0][]02W5}{[1][0][]source}{}{}{}}
\makeatother
\crefname{originalreference}{original reference}{original references}
\Crefname{originalreference}{Original reference}{Original references}
\makeatletter
\newlabel{lemma-F-zero-module-map}{{083L}{source}{}{sourcepointer-3}{}}
\newlabel{lemma-F-zero-module-map@cref}{{[originalreference][0][]083L}{[1][0][]source}{}{}{}}
\makeatother
\crefname{originalreference}{original reference}{original references}
\Crefname{originalreference}{Original reference}{Original references}
\makeatletter
\newlabel{lemma-finite-presentation-flat-along-fibre}{{05KW}{source}{}{sourcepointer-4}{}}
\newlabel{lemma-finite-presentation-flat-along-fibre@cref}{{[originalreference][0][]05KW}{[1][0][]source}{}{}{}}
\makeatother
\crefname{originalreference}{original reference}{original references}
\Crefname{originalreference}{Original reference}{Original references}
\makeatletter
\newlabel{lemma-finite-type-flat-pure-is-universal}{{05L6}{source}{}{sourcepointer-5}{}}
\newlabel{lemma-finite-type-flat-pure-is-universal@cref}{{[originalreference][0][]05L6}{[1][0][]source}{}{}{}}
\makeatother
\crefname{originalreference}{original reference}{original references}
\Crefname{originalreference}{Original reference}{Original references}
\makeatletter
\newlabel{lemma-iso-limits}{{083J}{source}{}{sourcepointer-6}{}}
\newlabel{lemma-iso-limits@cref}{{[originalreference][0][]083J}{[1][0][]source}{}{}{}}
\makeatother
\crefname{originalreference}{original reference}{original references}
\Crefname{originalreference}{Original reference}{Original references}
\makeatletter
\newlabel{lemma-iso-sheaf}{{05MJ}{source}{}{sourcepointer-7}{}}
\newlabel{lemma-iso-sheaf@cref}{{[originalreference][0][]05MJ}{[1][0][]source}{}{}{}}
\makeatother
\crefname{originalreference}{original reference}{original references}
\Crefname{originalreference}{Original reference}{Original references}
\makeatletter
\newlabel{lemma-pure-on-top}{{0CW3}{source}{}{sourcepointer-8}{}}
\newlabel{lemma-pure-on-top@cref}{{[originalreference][0][]0CW3}{[1][0][]source}{}{}{}}
\makeatother
\crefname{originalreference}{original reference}{original references}
\Crefname{originalreference}{Original reference}{Original references}
\makeatletter
\newlabel{lemma-universally-separating}{{05PD}{source}{}{sourcepointer-9}{}}
\newlabel{lemma-universally-separating@cref}{{[originalreference][0][]05PD}{[1][0][]source}{}{}{}}
\makeatother
\crefname{originallemma}{original lemma}{original lemmas}
\Crefname{originallemma}{Original lemma}{Original lemmas}
\makeatletter
\newlabel{morphisms-lemma-constructible-containing-open}{{05LW}{source}{}{sourcepointer-10}{}}
\newlabel{morphisms-lemma-constructible-containing-open@cref}{{[originallemma][0][]05LW}{[1][0][]source}{}{}{}}
\makeatother
\crefname{originalreference}{original reference}{original references}
\Crefname{originalreference}{Original reference}{Original references}
\makeatletter
\newlabel{situation-iso}{{05MH}{source}{}{sourcepointer-11}{}}
\newlabel{situation-iso@cref}{{[originalreference][0][]05MH}{[1][0][]source}{}{}{}}
\makeatother
\crefname{originallemma}{original lemma}{original lemmas}
\Crefname{originallemma}{Original lemma}{Original lemmas}
\makeatletter
\newlabel{spaces-divisors-lemma-relative-assassin-constructible}{{0CVV}{source}{}{sourcepointer-12}{}}
\newlabel{spaces-divisors-lemma-relative-assassin-constructible@cref}{{[originallemma][0][]0CVV}{[1][0][]source}{}{}{}}
\makeatother
\crefname{originallemma}{original lemma}{original lemmas}
\Crefname{originallemma}{Original lemma}{Original lemmas}
\makeatletter
\newlabel{spaces-lemma-representable-over-space}{{02WY}{source}{}{sourcepointer-13}{}}
\newlabel{spaces-lemma-representable-over-space@cref}{{[originallemma][0][]02WY}{[1][0][]source}{}{}{}}
\makeatother
\crefname{originalproposition}{original proposition}{original propositions}
\Crefname{originalproposition}{Original proposition}{Original propositions}
\makeatletter
\newlabel{spaces-limits-proposition-characterize-locally-finite-presentation}{{04AK}{source}{}{sourcepointer-14}{}}
\newlabel{spaces-limits-proposition-characterize-locally-finite-presentation@cref}{{[originalproposition][0][]04AK}{[1][0][]source}{}{}{}}
\makeatother
\title{The zero locus of a map into a pure finite presentation module is a closed immersion}
\author{The Stacks Project Authors}
\date{stacks-project:spaces-flat.tex:local-snapshot-20261012; GNU FDL 1.2 or later}
\begin{document}
\maketitle
\noindent\textit{Editorial note.} All delivered mathematical text is the author's own, copied from the pinned Stacks Project chapter file spaces-flat.tex; no statement or proof step was rewritten, reordered or removed.\par\medskip
\noindent\textit{Editorial note.} The counted proof is the author's complete argument: a reduction to scheme bases and an affine chart on which the module is projective, a constructibility and purity argument shrinking the base, and an application of the universally separating property\par\medskip
\noindent\textit{Editorial note.} The label lemma-F-zero-closed-pure names the counted statement, and the proof environment immediately following it in the source is delivered verbatim as its proof.\par\medskip
\noindent\textit{Editorial note.} The counted statement is carried by the author's own lemma environment in the chapter source.\par\medskip
\noindent\textit{Allowed standard prerequisites (editorial).}\begin{itemize}
\item Universal purity of a flat module of finite presentation that is pure relative to the base, and the purity-on-top criterion relating purity to the associated points of the fibres
\item The criterion reducing representability to representability over schemes, and etale locality of closed immersions and of descent data for sheaves
\item The existence of a chart over an affine etale neighbourhood on which the pullback of the module is projective and whose image contains the fibre, the constructibility of the relative assassin, and the criterion that a map is injective exactly when its pullback along a universally separating morphism is
\item The affine case for a projective module over the coordinate ring of an affine base
\end{itemize}\medskip
\ifdefined\hypertarget\hypertarget{sourcepointer-1}{}\fi
\noindent\textit{Source reference.} Original lemma 03I0: sources/stacks-chow/tags.txt, source lines 4586--4586. Prerequisite pointer invoked at proof line(s) 1733; it names a standard fact of a sibling chapter that the counted argument uses, and every step of the counted proof built on it is retained verbatim in the delivered body.\par\medskip
\ifdefined\hypertarget\hypertarget{sourcepointer-2}{}\fi
\noindent\textit{Source reference.} Original lemma 02W5: sources/stacks-chow/tags.txt, source lines 3792--3792. Prerequisite pointer invoked at proof line(s) 1734; it names a standard fact of a sibling chapter that the counted argument uses, and every step of the counted proof built on it is retained verbatim in the delivered body.\par\medskip
\ifdefined\hypertarget\hypertarget{sourcepointer-3}{}\fi
\noindent\textit{Source reference.} Original reference 083L: sources/stacks-spaces-flat/pinned-original-source.tex, source lines 1627--1627. Prerequisite pointer invoked at proof line(s) 1779; it names a standard fact of the same chapter that the counted argument uses, and every step of the counted proof built on it is retained verbatim in the delivered body.\par\medskip
\ifdefined\hypertarget\hypertarget{sourcepointer-4}{}\fi
\noindent\textit{Source reference.} Original reference 05KW: sources/stacks-spaces-flat/pinned-original-source.tex, source lines 956--956. Prerequisite pointer invoked at proof line(s) 1745; it names a standard fact of the same chapter that the counted argument uses, and every step of the counted proof built on it is retained verbatim in the delivered body.\par\medskip
\ifdefined\hypertarget\hypertarget{sourcepointer-5}{}\fi
\noindent\textit{Source reference.} Original reference 05L6: sources/stacks-spaces-flat/pinned-original-source.tex, source lines 1146--1146. Prerequisite pointer invoked at proof line(s) 1720; it names a standard fact of the same chapter that the counted argument uses, and every step of the counted proof built on it is retained verbatim in the delivered body.\par\medskip
\ifdefined\hypertarget\hypertarget{sourcepointer-6}{}\fi
\noindent\textit{Source reference.} Original reference 083J: sources/stacks-spaces-flat/pinned-original-source.tex, source lines 1335--1335. Prerequisite pointer invoked at proof line(s) 1784; it names a standard fact of the same chapter that the counted argument uses, and every step of the counted proof built on it is retained verbatim in the delivered body.\par\medskip
\ifdefined\hypertarget\hypertarget{sourcepointer-7}{}\fi
\noindent\textit{Source reference.} Original reference 05MJ: sources/stacks-spaces-flat/pinned-original-source.tex, source lines 1269--1269. Prerequisite pointer invoked at proof line(s) 1732; it names a standard fact of the same chapter that the counted argument uses, and every step of the counted proof built on it is retained verbatim in the delivered body.\par\medskip
\ifdefined\hypertarget\hypertarget{sourcepointer-8}{}\fi
\noindent\textit{Source reference.} Original reference 0CW3: sources/stacks-spaces-flat/pinned-original-source.tex, source lines 1050--1050. Prerequisite pointer invoked at proof line(s) 1764; it names a standard fact of the same chapter that the counted argument uses, and every step of the counted proof built on it is retained verbatim in the delivered body.\par\medskip
\ifdefined\hypertarget\hypertarget{sourcepointer-9}{}\fi
\noindent\textit{Source reference.} Original reference 05PD: sources/stacks-spaces-flat/pinned-original-source.tex, source lines 1163--1163. Prerequisite pointer invoked at proof line(s) 1772; it names a standard fact of the same chapter that the counted argument uses, and every step of the counted proof built on it is retained verbatim in the delivered body.\par\medskip
\ifdefined\hypertarget\hypertarget{sourcepointer-10}{}\fi
\noindent\textit{Source reference.} Original lemma 05LW: sources/stacks-chow/tags.txt, source lines 7250--7250. Prerequisite pointer invoked at proof line(s) 1766; it names a standard fact of a sibling chapter that the counted argument uses, and every step of the counted proof built on it is retained verbatim in the delivered body.\par\medskip
\ifdefined\hypertarget\hypertarget{sourcepointer-11}{}\fi
\noindent\textit{Source reference.} Original reference 05MH: sources/stacks-spaces-flat/pinned-original-source.tex, source lines 1233--1233. Prerequisite pointer invoked at proof line(s) 1708; it names a standard fact of the same chapter that the counted argument uses, and every step of the counted proof built on it is retained verbatim in the delivered body.\par\medskip
\ifdefined\hypertarget\hypertarget{sourcepointer-12}{}\fi
\noindent\textit{Source reference.} Original lemma 0CVV: sources/stacks-chow/tags.txt, source lines 16150--16150. Prerequisite pointer invoked at proof line(s) 1759; it names a standard fact of a sibling chapter that the counted argument uses, and every step of the counted proof built on it is retained verbatim in the delivered body.\par\medskip
\ifdefined\hypertarget\hypertarget{sourcepointer-13}{}\fi
\noindent\textit{Source reference.} Original lemma 02WY: sources/stacks-chow/tags.txt, source lines 3821--3821. Prerequisite pointer invoked at proof line(s) 1724; it names a standard fact of a sibling chapter that the counted argument uses, and every step of the counted proof built on it is retained verbatim in the delivered body.\par\medskip
\ifdefined\hypertarget\hypertarget{sourcepointer-14}{}\fi
\noindent\textit{Source reference.} Original proposition 04AK: sources/stacks-chow/tags.txt, source lines 5619--5619. Prerequisite pointer invoked at proof line(s) 1786; it names a standard fact of a sibling chapter that the counted argument uses, and every step of the counted proof built on it is retained verbatim in the delivered body.\par\medskip
\begin{lemma}
\label{lemma-F-zero-closed-pure}
In Situation \ref{situation-iso}. Assume
\begin{enumerate}
\item $f$ is of finite presentation, and
\item $\mathcal{G}$ is of finite presentation,
flat over $B$, and pure relative to $B$.
\end{enumerate}
Then $F_{zero}$ is an algebraic space and $F_{zero} \to B$
is a closed immersion. If $\mathcal{F}$ is of finite type, then
$F_{zero} \to B$ is of finite presentation.
\end{lemma}

\begin{proof}
By Lemma \ref{lemma-finite-type-flat-pure-is-universal}
the module $\mathcal{G}$ is universally pure relative to $B$.
In order to prove that $F_{zero}$ is an algebraic space,
it suffices to show that $F_{zero} \to B$ is representable, see
Spaces, Lemma \ref{spaces-lemma-representable-over-space}.
Let $B' \to B$ be a morphism where $B'$ is a scheme and let
$u' : \mathcal{F}' \to \mathcal{G}'$ be the pullback of $u$ to $X' = X_{B'}$.
Then the associated functor $F'_{zero}$ equals $F_{zero} \times_B B'$.
This reduces us to the case that $B$ is a scheme.

\medskip\noindent
Assume $B$ is a scheme. We will show that $F_{zero}$ is representable
by a closed subscheme of $B$. By Lemma \ref{lemma-iso-sheaf} and
Descent, Lemmas \ref{descent-lemma-closed-immersion} and
\ref{descent-lemma-descent-data-sheaves}
the question is local for the \'etale topology on $B$. Let $b \in B$.
We first replace $B$ by an affine neighbourhood of $b$.
Choose a diagram
$$
\xymatrix{
X \ar[d] & X' \ar[l]^g \ar[d] \\
B & B' \ar[l]
}
$$
and $b' \in B'$ mapping to $b \in B$
as in Lemma \ref{lemma-finite-presentation-flat-along-fibre}.
As we are working \'etale locally, we may replace
$B$ by $B'$ and assume that we have a diagram
$$
\xymatrix{
X \ar[rd] & & X' \ar[ll]^g \ar[ld] \\
& B
}
$$
with $B$ and $X'$ affine such that $\Gamma(X', g^*\mathcal{G})$
is a projective $\Gamma(B, \mathcal{O}_B)$-module and
$g(|X'|) \supset |X_b|$. Let $U \subset X$ be the open subspace
with $|U| = g(|X'|)$. By
Divisors on Spaces, Lemma
\ref{spaces-divisors-lemma-relative-assassin-constructible} the set
$$
E = \{t \in B : \text{Ass}_{X_t}(\mathcal{G}_t) \subset |U_t|\} =
\{t \in B : \text{Ass}_{X/B}(\mathcal{G}) \cap |X_t| \subset |U_t|\}
$$
is constructible in $B$. By Lemma \ref{lemma-pure-on-top} part (2)
we see that $E$ contains $\Spec(\mathcal{O}_{B, b})$. By
Morphisms, Lemma \ref{morphisms-lemma-constructible-containing-open}
we see that $E$ contains an open neighbourhood of $b$. Hence after
replacing $B$ by a smaller affine neighbourhood of $b$ we may assume that
$\text{Ass}_{X/B}(\mathcal{G}) \subset g(|X'|)$.

\medskip\noindent
From Lemma \ref{lemma-universally-separating}
it follows that $u : \mathcal{F} \to \mathcal{G}$ is injective if and only if
$g^*u : g^*\mathcal{F} \to g^*\mathcal{G}$ is injective, and the same remains
true after any base change. Hence we have reduced to the case where,
in addition to the assumptions in the theorem, $X \to B$ is a morphism of
affine schemes and $\Gamma(X, \mathcal{G})$ is a projective
$\Gamma(B, \mathcal{O}_B)$-module. This case follows immediately from
Lemma \ref{lemma-F-zero-module-map}.

\medskip\noindent
We still have to show that $F_{zero} \to B$ is of finite presentation if
$\mathcal{F}$ is of finite type. This follows from
Lemma \ref{lemma-iso-limits} combined with
Limits of Spaces, Proposition
\ref{spaces-limits-proposition-characterize-locally-finite-presentation}.
\end{proof}

\end{document}
